\chapter{Cofactors and Cramer's Rule}

In the previous chapter we learned what a determinant \emph{is}: a single number attached to a square
matrix that measures signed volume -- the factor by which the matrix scales areas in $\R^2$, volumes
in $\R^3$, and hyper-volumes beyond, together with a sign that records whether orientation is kept or
flipped. We learned the headline facts. The determinant is zero exactly when the matrix collapses
space (when its columns are dependent, when it is not invertible). It multiplies: $\det(\mat A \mat
B)=\det\mat A\,\det\mat B$. It is unchanged by the row operation that elimination uses most.

This chapter answers a different question. Not ``what does the determinant mean,'' but ``how is it
\emph{built}, and what can we build \emph{out} of it.'' Two classical structures grow directly from
the determinant: the \textbf{cofactor expansion}, which computes a big determinant by peeling it into
smaller ones, and from there the \textbf{adjugate}, which assembles those small determinants into an
explicit formula for the inverse. Cramer's rule then drops out as a one-line consequence, and it
carries a clean geometric reading: each unknown in a linear system is a ratio of two volumes.

A warning before we start, because it will shape how you read everything that follows. The formulas
in this chapter are some of the most elegant in linear algebra, but among the
least practical for numerical calculation. They are tools for
\emph{proving theorems} and for \emph{understanding algebraic structure}, not tools for crunching
numbers on large matrices. At the end of the chapter we compare their operation counts directly
against Gaussian elimination.

\begin{focusbox}
The whole chapter hangs on one identity, the \textbf{cofactor expansion}: a determinant equals a
sum, along any single row or column, of each entry times a $\pm$ sign times the determinant of the
smaller matrix you get by deleting that entry's row and column. Everything else -- the inverse
formula, Cramer's rule, the volume interpretation -- is a consequence of reading that one identity
the right way.
\end{focusbox}

\section{Peeling a Determinant Apart: Minors and Cofactors}

Let us start where the geometry is easiest to hold: a $2\times2$ matrix, whose determinant you
already know.
\[
\mat A=\begin{bmatrix} a & b \\ c & d \end{bmatrix},\qquad \det\mat A = ad-bc.
\]
Stare at $ad-bc$ and ask: where do the two terms come from? The $a$ multiplies the $d$ -- and $d$ is
exactly the lone entry left over when you cross out $a$'s row (the top) and $a$'s column (the left).
The $b$ multiplies $c$ -- and $c$ is what is left when you cross out $b$'s row and column. So the
$2\times2$ determinant is already a sum over the top row: take each top entry, multiply it by ``the
determinant of what remains,'' and attach an alternating sign $+,-$.

That little observation is the entire idea, repeated at every size. Let us name the pieces.

\begin{definition}[Minor]\label{ch22:def-minor}
Let $\mat A$ be an $n\times n$ matrix. The \textbf{minor} $M_{ij}$ is the determinant of the
$(n-1)\times(n-1)$ matrix obtained by deleting row $i$ and column $j$ of $\mat A$.
\end{definition}

\begin{definition}[Cofactor]\label{ch22:def-cofactor}
The \textbf{cofactor} $C_{ij}$ is the signed minor
\[
C_{ij} = (-1)^{\,i+j}\,M_{ij}.
\]
The sign $(-1)^{i+j}$ is $+1$ when $i+j$ is even and $-1$ when $i+j$ is odd.
\end{definition}

The sign pattern follows a checkerboard layout, always starting with $+$ in the top-left corner:
\[
\begin{bmatrix}
+ & - & + & \cdots \\
- & + & - & \cdots \\
+ & - & + & \cdots \\
\vdots & & & \ddots
\end{bmatrix}.
\]

\begin{intuitionbox}
A minor answers ``forget about entry $(i,j)$ entirely -- if I delete its whole row and column, how
much volume does the leftover matrix make?'' The cofactor is that same number with a sign glued on.
The sign is not decoration: it is bookkeeping for orientation, the same $\pm$ that distinguishes a
rotation from a reflection. When we add up entry $\times$ cofactor, the signs make the cross-terms
cancel in exactly the way that recovers signed volume.
\end{intuitionbox}

\begin{trapbox}
The checkerboard sign $(-1)^{i+j}$ belongs to the \emph{cofactor}, not the minor. A frequent error is
to compute the minors correctly, then forget the sign on the off-diagonal positions and get the wrong
total. Equally common: starting the checkerboard with a $-$ in the top-left. It always starts with
$+$, because $(-1)^{1+1}=(-1)^2=+1$.
\end{trapbox}

\begin{labbox}
\textbf{Pause and think:} What is the checkerboard sign for entry $(2,3)$? For entry $(3,1)$?
Count with your fingers: row 2, column 3 gives $(-1)^{2+3} = (-1)^5 = -1$. Row 3, column 1 gives
$(-1)^{3+1} = (-1)^4 = +1$. Whenever the sum of coordinates is odd, you get a minus; whenever even, a plus.
\end{labbox}

\subsection{The Cofactor Expansion Theorem}

Now the central claim. You can compute the determinant by marching along \emph{any} row or
\emph{any} column, multiplying each entry by its cofactor, and adding.

\begin{theorem}[Cofactor / Laplace expansion]\label{ch22:thm-expansion}
Let $\mat A$ be $n\times n$ with entries $a_{ij}$. For any fixed row $i$,
\[
\det\mat A \;=\; \sum_{j=1}^{n} a_{ij}\,C_{ij}
\;=\; \sum_{j=1}^{n} (-1)^{i+j}\,a_{ij}\,M_{ij},
\]
and for any fixed column $j$,
\[
\det\mat A \;=\; \sum_{i=1}^{n} a_{ij}\,C_{ij}
\;=\; \sum_{i=1}^{n} (-1)^{i+j}\,a_{ij}\,M_{ij}.
\]
\end{theorem}

The key word in that statement is \emph{any}. Expand along the top row, the bottom row, the
middle column -- you get the same number every time. We will not give the full combinatorial proof
from the permutation definition (that is a worthwhile but heavy bookkeeping exercise), but the
geometric and recursive heart of it is transparent, and we sketch it in the box below.

\begin{formalbox}
\textbf{Why expansion works (sketch).} The determinant can be \emph{defined} as the signed sum over
all permutations $\sigma$ of $\{1,\dots,n\}$,
\[
\det\mat A=\sum_{\sigma}\operatorname{sgn}(\sigma)\,a_{1\sigma(1)}a_{2\sigma(2)}\cdots a_{n\sigma(n)}.
\]
Fix row $i$. Every term in that giant sum uses \emph{exactly one} entry from row $i$, say
$a_{i\sigma(i)}$. Group the terms by which column $j$ that entry sits in. The terms with
$\sigma(i)=j$, after factoring out $a_{ij}$, are precisely a signed sum over permutations of the
remaining $n-1$ indices -- which is the determinant of the matrix with row $i$ and column $j$
removed. Tracking how $\operatorname{sgn}(\sigma)$ splits gives exactly the factor $(-1)^{i+j}$. So
the grouped sum is $\sum_j a_{ij}(-1)^{i+j}M_{ij}$. Because the original sum was symmetric in how we
chose to single out a row, the same argument works for any row or column.
\end{formalbox}

\subsection{A worked \texorpdfstring{$3\times3$}{3x3} expansion}

Let us actually do one, slowly, because the mechanics are where people stumble.
\[
\mat A=\begin{bmatrix} 2 & -1 & 3 \\ 0 & 4 & 5 \\ 1 & 0 & 6 \end{bmatrix}.
\]
We expand along the \textbf{first row}. The three top entries are $2,-1,3$, sitting in positions
$(1,1)$, $(1,2)$, $(1,3)$, with signs $+,-,+$.

\emph{First term} -- entry $2$, sign $+$. Delete row $1$, column $1$:
\[
M_{11}=\det\begin{bmatrix} 4 & 5 \\ 0 & 6 \end{bmatrix}=4\cdot6-5\cdot0=24,
\qquad +\,2\cdot 24 = 48.
\]
\emph{Second term} -- entry $-1$, sign $-$. Delete row $1$, column $2$:
\[
M_{12}=\det\begin{bmatrix} 0 & 5 \\ 1 & 6 \end{bmatrix}=0\cdot6-5\cdot1=-5,
\qquad -\,(-1)\cdot(-5) = -5.
\]
\emph{Third term} -- entry $3$, sign $+$. Delete row $1$, column $3$:
\[
M_{13}=\det\begin{bmatrix} 0 & 4 \\ 1 & 0 \end{bmatrix}=0\cdot0-4\cdot1=-4,
\qquad +\,3\cdot(-4) = -12.
\]
Add them: $\det\mat A = 48-5-12 = 31$.

Now the payoff of ``any row or column.'' That first column of $\mat A$ has a zero in it ($a_{21}=0$)
and a one ($a_{31}=1$), which makes the arithmetic lighter. Expand down \textbf{column 1} instead,
signs $+,-,+$:
\[
\det\mat A
= +\,2\det\begin{bmatrix}4&5\\0&6\end{bmatrix}
 -\,0\cdot(\cdots)
 +\,1\det\begin{bmatrix}-1&3\\4&5\end{bmatrix}
= 2(24) + 1\big((-1)(5)-3\cdot4\big)
= 48 + (-17) = 31.
\]
Same $31$, and the middle term vanished for free. This is the one genuine computational lesson of the
section, so we make it a rule.

\begin{focusbox}
\textbf{Expand along the row or column with the most zeros.} Each zero entry kills an entire cofactor
-- a whole smaller determinant you never have to compute. Choosing wisely can turn a painful
expansion into a short one.
\end{focusbox}

\begin{labbox}
Compute $\det\!\begin{bmatrix} 1 & 2 & 0 \\ 3 & -1 & 0 \\ 4 & 5 & 2 \end{bmatrix}$ by hand. The third
column has two zeros -- expand along it. You should find only one nonzero cofactor survives, and it is
$2$ times a $2\times2$ determinant. (Answer: $2\cdot(1\cdot(-1)-2\cdot3)=2\cdot(-7)=-14$.) Now redo
it expanding along the first row and confirm you get $-14$ again, with more work.
\end{labbox}

Here is a picture of the peeling process: one $3\times3$ determinant becomes three $2\times2$
determinants, each living on the entries that survive a deletion.

\begin{figure}[H]
\centering
\begin{tikzpicture}[scale=0.62, every node/.style={font=\small}]
  % big matrix
  \foreach \r in {0,1,2}{
    \foreach \c in {0,1,2}{
      \draw[gray!40] (\c,-\r) rectangle ++(1,-1);
    }
  }
  \node at (0.5,-0.5){$a_{11}$};
  \node at (1.5,-0.5){$a_{12}$};
  \node at (2.5,-0.5){$a_{13}$};
  \node at (0.5,-1.5){$a_{21}$};
  \node at (1.5,-1.5){$a_{22}$};
  \node at (2.5,-1.5){$a_{23}$};
  \node at (0.5,-2.5){$a_{31}$};
  \node at (1.5,-2.5){$a_{32}$};
  \node at (2.5,-2.5){$a_{33}$};
  % highlight top row
  \draw[vecblue,very thick] (0,0) rectangle (3,-1);
  \node[vecblue] at (1.5,0.5){expand along this row};

  % arrows to three minors
  \begin{scope}[xshift=6.5cm,yshift=0.2cm]
    \node[vecred, font=\small\bfseries] (t1) at (-1.1,-1) {$+\,a_{11}\cdot$};
    \draw[thick] (-0.1,0.1) -- (-0.1,-2.1);
    \draw[thick] (2.1,0.1) -- (2.1,-2.1);
    \draw[gray!50] (0,0) rectangle (1,-1);
    \draw[gray!50] (1,0) rectangle (2,-1);
    \draw[gray!50] (0,-1) rectangle (1,-2);
    \draw[gray!50] (1,-1) rectangle (2,-2);
    \node at (0.5,-0.5){$a_{22}$}; \node at (1.5,-0.5){$a_{23}$};
    \node at (0.5,-1.5){$a_{32}$}; \node at (1.5,-1.5){$a_{33}$};
  \end{scope}
  \begin{scope}[xshift=6.5cm,yshift=-3cm]
    \node[vecred, font=\small\bfseries] (t2) at (-1.1,-1) {$-\,a_{12}\cdot$};
    \draw[thick] (-0.1,0.1) -- (-0.1,-2.1);
    \draw[thick] (2.1,0.1) -- (2.1,-2.1);
    \draw[gray!50] (0,0) rectangle (1,-1);
    \draw[gray!50] (1,0) rectangle (2,-1);
    \draw[gray!50] (0,-1) rectangle (1,-2);
    \draw[gray!50] (1,-1) rectangle (2,-2);
    \node at (0.5,-0.5){$a_{21}$}; \node at (1.5,-0.5){$a_{23}$};
    \node at (0.5,-1.5){$a_{31}$}; \node at (1.5,-1.5){$a_{33}$};
  \end{scope}
  \begin{scope}[xshift=6.5cm,yshift=-6cm]
    \node[vecred, font=\small\bfseries] (t3) at (-1.1,-1) {$+\,a_{13}\cdot$};
    \draw[thick] (-0.1,0.1) -- (-0.1,-2.1);
    \draw[thick] (2.1,0.1) -- (2.1,-2.1);
    \draw[gray!50] (0,0) rectangle (1,-1);
    \draw[gray!50] (1,0) rectangle (2,-1);
    \draw[gray!50] (0,-1) rectangle (1,-2);
    \draw[gray!50] (1,-1) rectangle (2,-2);
    \node at (0.5,-0.5){$a_{21}$}; \node at (1.5,-0.5){$a_{22}$};
    \node at (0.5,-1.5){$a_{31}$}; \node at (1.5,-1.5){$a_{32}$};
  \end{scope}
  \draw[-{Stealth}, vecblue!70!black, thick] (3.1,-0.5) .. controls (4.4,-0.5) .. (t1.west);
  \draw[-{Stealth}, vecblue!70!black, thick] (3.1,-0.5) .. controls (4.4,-2.2) .. (t2.west);
  \draw[-{Stealth}, vecblue!70!black, thick] (3.1,-0.5) .. controls (4.4,-4.5) .. (t3.west);
\end{tikzpicture}
\caption{Cofactor expansion along the top row turns one $3\times3$ determinant into three signed
$2\times2$ minors. Each minor lives on the entries that survive deleting the highlighted entry's row
and column.}
\label{ch22:fig-peel}
\end{figure}

\begin{strangbox}
Strang introduces cofactors as the third of his ``three ways to compute a determinant'' (the others
being elimination to triangular form, and the big permutation formula). Reading this \emph{before}
Strang, hold onto the recursive picture in Figure~\ref{ch22:fig-peel}: a determinant of size $n$ is a
signed combination of determinants of size $n-1$. Reading it \emph{after} Strang, notice that his
``BIG formula'' (the permutation sum) and the cofactor formula are the same sum, merely grouped
differently -- grouping by the choice in one row is exactly what produces the cofactors.
\end{strangbox}

\subsection{Cofactors and the cost of computing}

It is tempting to think cofactor expansion is \emph{the} way to find a determinant. For tiny matrices
it is fine. But count the work. To expand an $n\times n$ determinant you compute $n$ cofactors, each
an $(n-1)\times(n-1)$ determinant; each of those spawns $n-1$ more, and so on. The number of
multiplications grows like $n!$ -- factorial. For $n=10$ that is over three million products; for
$n=20$ it is more than the number of seconds in the age of the universe. Elimination, by contrast,
reaches a triangular form (whose determinant is just the product of the pivots) in roughly $n^3/3$
operations. For $n=20$ that is a few thousand. We will return to this contrast in
Section~\ref{ch22:sec-when}, but we emphasize this distinction: \textbf{cofactor expansion is a reasoning tool,
not a computing tool.}

The difference in efficiency is asymptotic. Figure~\ref{ch22:fig-cost} compares the two operation
counts on a logarithmic vertical axis, illustrating how quickly the factorial curve $n!$ outpaces
the cubic curve $n^3/3$.

\begin{figure}[H]
\centering
\begin{tikzpicture}
\begin{axis}[
    width=12cm, height=7cm,
    ymode=log,
    xlabel={matrix size $n$},
    ylabel={multiplications (log scale)},
    xmin=1, xmax=15, ymin=1, ymax=1e13,
    grid=both, grid style={gray!18},
    legend pos=north west,
    legend cell align=left,
    legend style={fill=white, fill opacity=0.95, draw=gray!40, text opacity=1},
    every axis plot/.append style={very thick},
    tick label style={font=\small}, label style={font=\small},
]
  % factorial cost ~ n!
  \addplot[vecred, mark=*, mark size=1.4pt] coordinates {
    (1,1)(2,2)(3,6)(4,24)(5,120)(6,720)(7,5040)(8,40320)(9,362880)
    (10,3628800)(11,39916800)(12,479001600)(13,6227020800)(14,87178291200)(15,1307674368000)
  };
  \addlegendentry{cofactor expansion $\sim n!$}
  % elimination ~ n^3/3
  \addplot[sourceblue, mark=square*, mark size=1.4pt] coordinates {
    (1,1)(2,3)(3,9)(4,21)(5,42)(6,72)(7,114)(8,171)(9,243)
    (10,333)(11,444)(12,576)(13,732)(14,915)(15,1125)
  };
  \addlegendentry{elimination $\sim n^3/3$}
\end{axis}
\end{tikzpicture}
\caption{The computational cost of evaluating a determinant. On a logarithmic scale, the cofactor
operation count $n!$ (red) grows rapidly, while Gaussian elimination's $n^3/3$ count (blue) remains
manageable. By $n=15$, cofactor expansion requires over a trillion multiplications, compared to
approximately one thousand for elimination.}
\label{ch22:fig-cost}
\end{figure}

\section{The Determinant of Special Matrices}

Before deriving the inverse formula, we apply the cofactor expansion to several special
families of matrices whose determinants can be evaluated immediately by inspection.

\subsection{Triangular matrices: just multiply the diagonal}

An \textbf{upper triangular} matrix has all zeros below its main diagonal, so each row only involves its
own variable and later ones. A \textbf{lower triangular} matrix has all zeros above its main diagonal.

\begin{theorem}[Determinant of a triangular matrix]\label{ch22:thm-triangular}
If $\mat A$ is upper triangular or lower triangular, then
\[
\det\mat A = a_{11}\,a_{22}\,\cdots\,a_{nn},
\]
the product of its diagonal entries.
\end{theorem}

\begin{proof}
Take $\mat A$ upper triangular, so every entry below the diagonal is zero. Expand along the
\emph{first column}. Only $a_{11}$ is nonzero there, so
\[
\det\mat A = a_{11}\,C_{11}=a_{11}\,M_{11},
\]
and the minor $M_{11}$ is the determinant of the $(n-1)\times(n-1)$ matrix you get by deleting the
first row and column -- which is \emph{still upper triangular}, with diagonal $a_{22},\dots,a_{nn}$.
Expanding it along \emph{its} first column peels off $a_{22}$, and so on by induction until a single
diagonal entry remains. Multiplying the peeled factors gives $a_{11}a_{22}\cdots a_{nn}$. The
lower-triangular case is identical, expanding along the first row each time.
\end{proof}

\begin{intuitionbox}
Geometrically a triangular matrix stretches space one axis at a time without ever mixing a later axis
back into an earlier one. The volume scaling is then just the product of the per-axis stretches --
exactly the diagonal entries. A zero on the diagonal means one axis got flattened, the volume is $0$,
and the matrix is singular.
\end{intuitionbox}

Figure~\ref{ch22:fig-triangular-shear} shows this geometrically in the plane. The diagonal entries
$a_{11}$ and $a_{22}$ determine the base length and the vertical height of the transformed box. The entry
above the diagonal, $a_{12}$, acts as a pure horizontal shear: it slides the top edge sideways parallel to
the base. By Cavalieri's principle ($\text{base}\times\text{height}$), sliding the ceiling sideways changes
the tilt of the walls but leaves the enclosed area untouched. The shear entry $a_{12}$ has no say in the
determinant.

\begin{figure}[H]
\centering
\begin{tikzpicture}[scale=1.05, >=Stealth, font=\small]
  \begin{scope}[shift={(0,0)}]
    \draw[->, spacegray] (-0.3,0) -- (2.2,0) node[right, font=\footnotesize] {$x$};
    \draw[->, spacegray] (0,-0.3) -- (0,2.2) node[above, font=\footnotesize] {$y$};
    \fill[sourceblue!15] (0,0) -- (1.4,0) -- (1.4,1.4) -- (0,1.4) -- cycle;
    \draw[sourceblue!60, dashed] (0,0) -- (1.4,0) -- (1.4,1.4) -- (0,1.4) -- cycle;
    \draw[->, vecblue, thick] (0,0) -- (1.4,0) node[midway, below=2pt] {\footnotesize $\vec{e}_1$};
    \draw[->, targetgreen, thick] (0,0) -- (0,1.4) node[midway, left=2pt] {\footnotesize $\vec{e}_2$};
    \node[sourceblue!90!black, font=\footnotesize\bfseries] at (0.7,0.7) {Area $= 1$};
    \node[font=\footnotesize, spacegray] at (0.7,-0.65) {Unit square};
  \end{scope}

  \draw[->, very thick, spacegray] (2.4,0.9) -- (3.8,0.9)
    node[midway, above=2pt, font=\normalsize\bfseries, black] {$\mat A$};

  \begin{scope}[shift={(4.5,0)}]
    \def\aa{2.2}
    \def\shear{1.0}
    \def\h{1.5}
    \draw[->, spacegray] (-0.3,0) -- (3.8,0) node[right, font=\footnotesize] {$x$};
    \draw[->, spacegray] (0,-0.3) -- (0,2.2) node[above, font=\footnotesize] {$y$};
    \draw[gray!35, dashed] (0,0) -- (\aa,0) -- (\aa,\h) -- (0,\h) -- cycle;
    \fill[targetgreen!18] (0,0) -- (\aa,0) -- (\aa+\shear,\h) -- (\shear,\h) -- cycle;
    \draw[targetgreen!70, thick] (0,0) -- (\aa,0) -- (\aa+\shear,\h) -- (\shear,\h) -- cycle;
    \draw[->, vecblue, very thick] (0,0) -- (\aa,0)
      node[midway, above=1pt] {\footnotesize $\mat A\vec{e}_1$};
    \draw[->, targetgreen, very thick] (0,0) -- (\shear,\h)
      node[midway, left=2pt] {\footnotesize $\mat A\vec{e}_2$};
    \draw[<->, trapred, thick] (0,-0.35) -- (\aa,-0.35)
      node[midway, below=1pt, font=\footnotesize] {$\text{base} = a_{11}$};
    \draw[<->, trapred, thick] (\aa+\shear+0.3, 0) -- (\aa+\shear+0.3, \h)
      node[midway, right=1pt, font=\footnotesize] {$\text{height} = a_{22}$};
    \draw[gray!50, dotted] (\aa,0) -- (\aa+\shear+0.35,0);
    \draw[gray!50, dotted] (\aa+\shear,\h) -- (\aa+\shear+0.35,\h);
    \draw[->, quotientorange, thick] (0,\h) -- (\shear,\h)
      node[midway, above=2pt, font=\footnotesize] {shear $a_{12}$};
    \node[targetgreen!90!black, font=\footnotesize\bfseries, fill=white, inner sep=1.5pt, rounded corners=2pt]
      at (1.6, 0.75) {Area $= a_{11}a_{22}$};
    \node[font=\footnotesize, spacegray] at (1.6,-0.85) {Sheared parallelogram};
  \end{scope}
\end{tikzpicture}
\caption{Why upper-triangular determinants ignore the entries above the diagonal. The diagonal entries
$a_{11}$ and $a_{22}$ fix the base length and perpendicular height. The entry above the diagonal, $a_{12}$,
acts as a horizontal shear, sliding the top edge sideways parallel to the base. By Cavalieri's principle,
sliding parallel to the base preserves area: $\text{Area} = \text{base}\times\text{height} = a_{11}a_{22}$,
completely unaffected by $a_{12}$.}
\label{ch22:fig-triangular-shear}
\end{figure}

This is precisely \emph{why} elimination computes determinants efficiently: row operations of the
``subtract a multiple of one row from another'' type do not change the determinant, and they drive any
matrix to triangular form, where the answer is a single product down the diagonal. Cofactors explain
\emph{why} the triangular form is so easy; elimination delivers you \emph{to} it cheaply.

\subsection{A worked diagonal and triangular check}

\[
\det\begin{bmatrix} 3 & 0 & 0 \\ 0 & -2 & 0 \\ 0 & 0 & 5 \end{bmatrix}=3\cdot(-2)\cdot5=-30,
\qquad
\det\begin{bmatrix} 4 & 7 & -1 \\ 0 & 2 & 9 \\ 0 & 0 & 6 \end{bmatrix}=4\cdot2\cdot6=48.
\]
The upper-triangular determinant ignores the entries $7,-1,9$ above the diagonal entirely. They
contribute nothing: as established in Theorem~\ref{ch22:thm-triangular}, expanding down successive
columns eliminates all above-diagonal contributions.

\subsection{A few more determinants you can read off}

\begin{itemize}
\item \textbf{Identity:} $\det\mat I = 1$. (Triangular with all $1$'s.) It scales no volume.
\item \textbf{A row or column of zeros:} $\det = 0$. Expand along that row; every term has a zero
  factor. Geometrically the matrix flattens space.
\item \textbf{Two equal rows (or columns):} $\det=0$. The columns are dependent, so volume collapses.
\item \textbf{Scaling the whole matrix:} $\det(c\mat A)=c^n\det\mat A$ for $n\times n$. Each of the
  $n$ axes is stretched by $c$, so volume scales by $c^n$ -- not by $c$. This trips people up
  constantly.
\item \textbf{Transpose:} $\det(\mat A\T)=\det\mat A$. The transpose $\mat A\T$ flips the matrix across its main diagonal, turning rows into columns. This equality is exactly why ``any row or column'' works
  in Theorem~\ref{ch22:thm-expansion}: the rows of $\mat A$ are the columns of $\mat A\T$.
\end{itemize}

\begin{trapbox}
$\det(c\mat A)=c^n\det\mat A$, not $c\det\mat A$. If you double every entry of a $3\times3$ matrix you
multiply its determinant by $2^3=8$, because you doubled all three edge-lengths of the box. Likewise
$\det(\mat A+\mat B)$ is \emph{not} $\det\mat A+\det\mat B$ in general -- the determinant is nonlinear
with respect to matrix addition, even though it is linear in each row \emph{separately}.
\end{trapbox}

\begin{labbox}
\textbf{Pause and think:} If $\mat A$ is a $4\times4$ matrix with $\det\mat A = 5$, what is $\det(-\mat A)$?
What is $\det(2\mat A)$? Before writing anything down, picture four orthogonal edges being scaled by $-1$
(four sign flips make $+1$, so $(-1)^4\cdot 5 = 5$) or by $2$ (each of the four dimensions doubles, scaling
the 4-dimensional hypervolume by $2^4 = 16$, giving $16\cdot 5 = 80$).
\end{labbox}

\section{The Adjugate and the Inverse Formula}

We now turn the cofactors into something structural. Collect every cofactor $C_{ij}$ into a matrix,
then \emph{transpose} it. That transposed grid is the object this section is built around.

\begin{definition}[Adjugate]\label{ch22:def-adj}
The \textbf{cofactor matrix} of $\mat A$ is the $n\times n$ matrix $\mat C$ whose $(i,j)$ entry is the
cofactor $C_{ij}$. The \textbf{adjugate} (or classical adjoint) is its transpose:
\[
\operatorname{adj}(\mat A) = \mat C\T,
\qquad\text{so}\qquad
\big(\operatorname{adj}\mat A\big)_{ij} = C_{ji}.
\]
\end{definition}

The name comes from the Latin \emph{adjugare} (``to yoke together to assist''), because it is yoked
to $\mat A$ to recover the identity matrix when multiplied. The transpose is easy to forget and it matters:
the $(i,j)$ entry of the adjugate is the cofactor $C_{ji}$ with the indices \emph{swapped}. Now the key identity.

\begin{theorem}[Adjugate identity]\label{ch22:thm-adj}
For any square $\mat A$,
\[
\mat A\,\operatorname{adj}(\mat A) = \operatorname{adj}(\mat A)\,\mat A = (\det\mat A)\,\mat I.
\]
Consequently, if $\det\mat A\neq 0$,
\[
\boxed{\;\mat A\inv = \frac{1}{\det\mat A}\,\operatorname{adj}(\mat A).\;}
\]
\end{theorem}

This formula provides a closed-form expression for the inverse of \emph{any}
invertible matrix without performing row reduction -- expressing each entry as a ratio of
sub-determinants over the determinant. The proof reveals the algebraic mechanism behind the formula.

\subsection{Why the adjugate identity holds}

Look at the $(i,i)$ diagonal entry of the product $\mat A\,\operatorname{adj}(\mat A)$. By the rule
for matrix multiplication it is
\[
\big(\mat A\,\operatorname{adj}\mat A\big)_{ii}
=\sum_{k} a_{ik}\big(\operatorname{adj}\mat A\big)_{ki}
=\sum_{k} a_{ik}\,C_{ik}.
\]
But $\sum_k a_{ik}C_{ik}$ is exactly the cofactor expansion of $\det\mat A$ along row $i$
(Theorem~\ref{ch22:thm-expansion})! So every diagonal entry equals $\det\mat A$. That handles the
diagonal.

Now the off-diagonal entry $(i,j)$ with $i\neq j$:
\[
\big(\mat A\,\operatorname{adj}\mat A\big)_{ij}
=\sum_{k} a_{ik}\big(\operatorname{adj}\mat A\big)_{kj}
=\sum_{k} a_{ik}\,C_{jk}.
\]
This sum pairs the entries of row $i$ with the cofactors of row $j$. Here is the clever observation:
$\sum_k a_{ik}C_{jk}$ is the cofactor expansion (along row $j$) of a \emph{modified} matrix -- the
matrix $\mat A$ with row $j$ \emph{replaced by a copy of row $i$}. Why? Because the cofactors $C_{jk}$
of row $j$ depend only on the \emph{other} rows; they never see row $j$ itself. So feeding them the
entries of row $i$ computes the determinant of a matrix that has row $i$ sitting in both position $i$
and position $j$. That matrix has two identical rows, so its determinant is $0$.

Therefore the product is $\det\mat A$ on the diagonal and $0$ everywhere else -- precisely $(\det\mat
A)\mat I$. Dividing by $\det\mat A$ (when it is nonzero) gives the inverse formula.

\begin{figure}[H]
\centering
\begin{tikzpicture}[scale=0.95, >=Stealth, font=\small]
  % Lane 1: Diagonal Entry (i,i)
  \begin{scope}[shift={(0,3.3)}]
    \node[anchor=west, font=\footnotesize\bfseries, text=vecblue!80!black] at (-0.2, 2.1)
      {Diagonal $(i,i)$: Row $i$ of $\mat{A}$ $\cdot$ Cofactors of Row $i$};
    % Box representation
    \draw[->, gray!50] (-0.3,0) -- (3.2,0);
    \draw[->, gray!50] (0,-0.3) -- (0,1.4);
    \fill[vecblue!18] (0,0) -- (2.2,0.3) -- (2.8,1.4) -- (0.6,1.1) -- cycle;
    \draw[vecblue, thick, ->] (0,0) -- (2.2,0.3) node[midway, below right=1pt] {\footnotesize $\vec{r}_i$};
    \draw[targetgreen, thick, ->] (0,0) -- (0.6,1.1) node[midway, above left=1pt] {\footnotesize other rows};
    \draw[vecblue!70, thick] (2.2,0.3) -- (2.8,1.4) -- (0.6,1.1);
    \node[fill=white, inner sep=1.5pt, rounded corners=2pt, font=\footnotesize\bfseries, text=vecblue!90!black] at (1.4,0.7) {Area $= \det\mat{A}$};
    \node[anchor=west, font=\footnotesize, align=left] at (5.0, 0.7)
      {$\displaystyle \sum_k a_{ik} C_{ik} = \det\mat{A}$\\[4pt]
       Standard cofactor expansion along row $i$.\\[2pt]
       Edges span full non-zero volume.};
  \end{scope}

  % Divider
  \draw[gray!30, dashed] (-0.2, 2.7) -- (11.8, 2.7);

  % Lane 2: Off-diagonal Entry (i,j)
  \begin{scope}[shift={(0,0)}]
    \node[anchor=west, font=\footnotesize\bfseries, text=trapred!80!black] at (-0.2, 2.1)
      {Off-Diagonal ($i \neq j$): Row $i$ of $\mat{A}$ $\cdot$ Cofactors of Row $j$};
    % Collapsed box
    \draw[->, gray!50] (-0.3,0) -- (3.2,0);
    \draw[->, gray!50] (0,-0.3) -- (0,1.4);
    % Row i twice:
    \draw[trapred, line width=1.4pt, ->] (0,0) -- (2.4,0.35) node[right=2pt] {\footnotesize $\vec{r}_i$ (row $i$)};
    \draw[quotientorange, line width=1.0pt, dashed, ->] (0,0.1) -- (2.4,0.5) node[above right=2pt] {\footnotesize $\vec{r}_i$ (row $j$)};
    \node[fill=white, inner sep=1.5pt, rounded corners=2pt, font=\footnotesize\bfseries, text=trapred!90!black] at (1.4,1.1) {Area $= 0$ (Squashed!)};
    \node[anchor=west, font=\footnotesize, align=left] at (5.0, 0.7)
      {$\displaystyle \sum_k a_{ik} C_{jk} = \det(\mat{A}\text{ with row }j \leftarrow \text{row }i) = 0$\\[4pt]
       Two identical rows $\vec{r}_i$ and $\vec{r}_i$.\\[2pt]
       Coinciding edges crush the box to zero volume.};
  \end{scope}
\end{tikzpicture}
\caption{Why the adjugate identity $\mat{A}\,\operatorname{adj}(\mat{A}) = (\det\mat{A})\,\mat{I}$ holds.
On the diagonal (top), dotting row $i$ with its own cofactors computes the full volume $\det\mat{A}$.
Off the diagonal (bottom), dotting row $i$ with the cofactors of row $j$ computes the volume of a box
where edge $j$ is replaced by edge $i$. Because two edges coincide, the box has zero width and collapses
to volume $0$.}
\label{ch22:fig-adjugate-geometry}
\end{figure}

\begin{aestheticbox}
The off-diagonal entries vanish because of a deliberate geometric cancellation: pairing one row's
entries with another row's cofactors computes the determinant of a box with two identical edges,
which has zero volume. That single observation -- duplicate rows mean zero volume -- makes the
inverse formula's off-diagonal entries disappear.
\end{aestheticbox}

\subsection{The \texorpdfstring{$2\times2$}{2x2} inverse, finally explained}

You have surely memorized the $2\times2$ inverse: ``swap the diagonal, negate the off-diagonal, divide
by the determinant.'' The adjugate formula is \emph{where that recipe comes from}. For
\[
\mat A=\begin{bmatrix} a & b \\ c & d \end{bmatrix},
\]
the four cofactors are single entries with signs: $C_{11}=+d$, $C_{12}=-c$, $C_{21}=-b$, $C_{22}=+a$.
The cofactor matrix is $\begin{bmatrix} d & -c \\ -b & a\end{bmatrix}$, and the adjugate is its
transpose:
\[
\operatorname{adj}\mat A=\begin{bmatrix} d & -b \\ -c & a\end{bmatrix},
\qquad
\mat A\inv=\frac{1}{ad-bc}\begin{bmatrix} d & -b \\ -c & a\end{bmatrix}.
\]
There it is. The mysterious swap-and-negate is the cofactor pattern $+,-,-,+$ followed by a transpose.

\begin{labbox}
Verify the adjugate identity by hand on $\mat A=\begin{bmatrix}2&1\\5&3\end{bmatrix}$. Here $\det\mat
A=1$. Form $\operatorname{adj}\mat A=\begin{bmatrix}3&-1\\-5&2\end{bmatrix}$ and multiply: $\mat
A\,\operatorname{adj}\mat A=\begin{bmatrix}2&1\\5&3\end{bmatrix}\begin{bmatrix}3&-1\\-5&2\end{bmatrix}
=\begin{bmatrix}1&0\\0&1\end{bmatrix}=(\det\mat A)\mat I.$ Because $\det\mat A=1$, the adjugate
\emph{is} the inverse here.
\end{labbox}

\subsection{A full \texorpdfstring{$3\times3$}{3x3} inverse by cofactors}

Let us do one $3\times3$ inverse from scratch, to see both the elegance and the labor. Take
\[
\mat A=\begin{bmatrix} 1 & 2 & 3 \\ 0 & 1 & 4 \\ 5 & 6 & 0 \end{bmatrix}.
\]
First the determinant; expand along column 1 (it has a zero):
\[
\det\mat A
= 1\det\begin{bmatrix}1&4\\6&0\end{bmatrix}
- 0
+ 5\det\begin{bmatrix}2&3\\1&4\end{bmatrix}
= 1(0-24) + 5(8-3) = -24 + 25 = 1.
\]
A determinant of $1$ -- convenient, because the inverse will equal the adjugate exactly. Now compute
all nine cofactors. Remember the sign checkerboard.
\[
\begin{aligned}
C_{11}&=+\det\begin{bmatrix}1&4\\6&0\end{bmatrix}=0-24=-24,
&C_{12}&=-\det\begin{bmatrix}0&4\\5&0\end{bmatrix}=-(0-20)=20,\\[2pt]
C_{13}&=+\det\begin{bmatrix}0&1\\5&6\end{bmatrix}=0-5=-5,
&C_{21}&=-\det\begin{bmatrix}2&3\\6&0\end{bmatrix}=-(0-18)=18,\\[2pt]
C_{22}&=+\det\begin{bmatrix}1&3\\5&0\end{bmatrix}=0-15=-15,
&C_{23}&=-\det\begin{bmatrix}1&2\\5&6\end{bmatrix}=-(6-10)=4,\\[2pt]
C_{31}&=+\det\begin{bmatrix}2&3\\1&4\end{bmatrix}=8-3=5,
&C_{32}&=-\det\begin{bmatrix}1&3\\0&4\end{bmatrix}=-(4-0)=-4,\\[2pt]
C_{33}&=+\det\begin{bmatrix}1&2\\0&1\end{bmatrix}=1-0=1.
\end{aligned}
\]
The cofactor matrix and its transpose (the adjugate):
\[
\mat C=\begin{bmatrix} -24 & 20 & -5 \\ 18 & -15 & 4 \\ 5 & -4 & 1\end{bmatrix},
\qquad
\operatorname{adj}\mat A=\mat C\T=\begin{bmatrix} -24 & 18 & 5 \\ 20 & -15 & -4 \\ -5 & 4 & 1\end{bmatrix}.
\]
Since $\det\mat A=1$,
\[
\mat A\inv=\operatorname{adj}\mat A=\begin{bmatrix} -24 & 18 & 5 \\ 20 & -15 & -4 \\ -5 & 4 & 1\end{bmatrix}.
\]
We can verify $\mat A\mat A\inv=\mat I$ directly. Notice the computational effort: nine $2\times2$
determinants plus the determinant of $\mat A$, all for a single inverse. Elimination produces the
same inverse with substantially fewer operations.

\begin{strangbox}
Strang is blunt about this formula, and rightly so: he calls the cofactor inverse a formula to
\emph{know} but rarely to \emph{use}. The reason to learn it is that it gives the inverse
\emph{symbolically} -- every entry of $\mat A\inv$ is a polynomial in the entries of $\mat A$ divided
by $\det\mat A$. That symbolic form is what lets you \emph{prove} general facts: that the inverse
entries are continuous in $\mat A$, that the inverse of an integer matrix with $\det=\pm1$ is again an
integer matrix, and more. Numbers come from elimination; \emph{theorems} often come from the
adjugate.
\end{strangbox}

\begin{connectionbox}
The adjugate gives a clean window into \emph{when integer matrices have integer inverses}, a question
at the heart of the group $\mathrm{GL}_n(\Z)$ and of lattices and cryptography. Because every entry of
$\operatorname{adj}\mat A$ is a sum of products of integer entries, it is an integer matrix whenever
$\mat A$ is. So $\mat A\inv=\frac{1}{\det\mat A}\operatorname{adj}\mat A$ is an integer matrix exactly
when $\det\mat A$ divides every adjugate entry -- and in particular always when $\det\mat A=\pm1$.
Matrices with $\det=\pm1$ are called \emph{unimodular}; they are precisely the invertible integer
matrices whose inverses are also integer matrices, and they form a group. This is why $\det=\pm1$
shows up everywhere from lattice-based cryptography to the change-of-basis matrices that preserve a
grid of points.
\end{connectionbox}

\section{Cramer's Rule}

Now we harvest the formula we have been building toward. Cramer's rule solves a square linear system
$\mat A\vec x=\vec b$ by writing \emph{each unknown} as a ratio of two determinants. No back
substitution, no augmented matrix -- just determinants.

\begin{theorem}[Cramer's rule]\label{ch22:thm-cramer}
Let $\mat A$ be $n\times n$ with $\det\mat A\neq 0$, and consider $\mat A\vec x=\vec b$. Let $\mat
A_i(\vec b)$ denote the matrix $\mat A$ with its $i$-th column replaced by $\vec b$. Then the unique
solution has entries
\[
x_i=\frac{\det \mat A_i(\vec b)}{\det\mat A},\qquad i=1,\dots,n.
\]
\end{theorem}

The single new idea in the statement is the matrix $\mat A_i(\vec b)$: take $\mat A$, throw away its
$i$-th column, and slot $\vec b$ into the gap. Every other column stays put. Figure~\ref{ch22:fig-replace}
shows the surgery for a $3\times3$ matrix, where the second column is the one replaced.

\begin{figure}[H]
\centering
\begin{tikzpicture}[scale=0.6, every node/.style={font=\small}]
  % original matrix A
  \foreach \c in {0,1,2}{\draw[gray!50] (\c,0) rectangle ++(1,-3);}
  \node at (0.5,-1.5){$\vec a_1$};
  \node[trapred] at (1.5,-1.5){$\vec a_2$};
  \node at (2.5,-1.5){$\vec a_3$};
  \node at (1.5,0.6){$\mat A$};
  % arrow
  \draw[-{Stealth},thick] (3.4,-1.5) -- (5.1,-1.5) node[midway,above]{replace col 2};
  % A_2(b)
  \begin{scope}[xshift=5.5cm]
    \foreach \c in {0,1,2}{\draw[gray!50] (\c,0) rectangle ++(1,-3);}
    \fill[vecred!18] (1,0) rectangle (2,-3);
    \node at (0.5,-1.5){$\vec a_1$};
    \node[vecred] at (1.5,-1.5){$\vec b$};
    \node at (2.5,-1.5){$\vec a_3$};
    \node at (1.5,0.6){$\mat A_2(\vec b)$};
  \end{scope}
\end{tikzpicture}
\caption{The matrix $\mat A_i(\vec b)$ is $\mat A$ with its $i$-th column overwritten by $\vec b$ (here
$i=2$). Cramer's rule says $x_i$ is the ratio $\det\mat A_i(\vec b)/\det\mat A$ of this matrix's
determinant to the original's.}
\label{ch22:fig-replace}
\end{figure}

\subsection{Where it comes from}

The cleanest derivation rides on the adjugate. From $\mat A\vec x=\vec b$ and $\mat A\inv=\frac{1}{\det
\mat A}\operatorname{adj}\mat A$ we get
\[
\vec x=\mat A\inv\vec b=\frac{1}{\det\mat A}\,\operatorname{adj}(\mat A)\,\vec b.
\]
The $i$-th entry of $\operatorname{adj}(\mat A)\,\vec b$ is $\sum_k\big(\operatorname{adj}\mat
A\big)_{ik}b_k=\sum_k C_{ki}\,b_k$. But $\sum_k b_k C_{ki}$ is exactly the cofactor expansion, down
column $i$, of the matrix whose $i$-th column is $\vec b$ and whose other columns are those of $\mat
A$ -- that is, $\det\mat A_i(\vec b)$. Hence $x_i=\det\mat A_i(\vec b)/\det\mat A$. Cramer's rule is
the adjugate inverse, read one coordinate at a time.

\begin{labbox}
\textbf{Pause and think:} Look at the formula $x_i = \det\mat A_i(\vec b)/\det\mat A$. What happens if $\vec b$
is already the $i$-th column of $\mat A$? Then $\mat A_i(\vec b)$ is simply $\mat A$, so $x_i = \det\mat A/\det\mat A = 1$.
And what happens to any other unknown $x_j$ ($j\neq i$)? In $\mat A_j(\vec b)$, column $j$ has been replaced by $\vec b$ (which is column $i$), so that matrix has two identical columns and its determinant is zero! Thus $x_j = 0$. In one line, Cramer's rule confirms that the $i$-th column vector has coordinates $(0,\dots,1,\dots,0)$.
\end{labbox}

\subsection{A worked \texorpdfstring{$2\times2$}{2x2} system}

Solve
\[
\begin{bmatrix} 2 & 1 \\ 1 & 3 \end{bmatrix}\begin{bmatrix} x_1 \\ x_2\end{bmatrix}
=\begin{bmatrix} 5 \\ 6\end{bmatrix}.
\]
First $\det\mat A=2\cdot3-1\cdot1=5$. For $x_1$, replace column 1 by $\vec b=\colvec{5\\6}$:
\[
\det\mat A_1(\vec b)=\det\begin{bmatrix} 5 & 1 \\ 6 & 3\end{bmatrix}=15-6=9,
\qquad x_1=\frac{9}{5}.
\]
For $x_2$, replace column 2 by $\vec b$:
\[
\det\mat A_2(\vec b)=\det\begin{bmatrix} 2 & 5 \\ 1 & 6\end{bmatrix}=12-5=7,
\qquad x_2=\frac{7}{5}.
\]
Check: $2\cdot\frac95+\frac75=\frac{18+7}{5}=5$ and $\frac95+3\cdot\frac75=\frac{9+21}{5}=6$. Both
equations hold.

\subsection{A worked \texorpdfstring{$3\times3$}{3x3} system}

Solve $\mat A\vec x=\vec b$ with
\[
\mat A=\begin{bmatrix} 1 & 0 & 2 \\ -1 & 3 & 1 \\ 2 & 1 & 1 \end{bmatrix},\qquad
\vec b=\begin{bmatrix} 1 \\ 2 \\ 5\end{bmatrix}.
\]
Determinant (expand along row 1):
\[
\det\mat A=1\det\begin{bmatrix}3&1\\1&1\end{bmatrix}-0+2\det\begin{bmatrix}-1&3\\2&1\end{bmatrix}
=1(3-1)+2(-1-6)=2-14=-12.
\]
Now three modified determinants. Replace column 1 by $\vec b$:
\[
\det\mat A_1=\det\begin{bmatrix} 1 & 0 & 2 \\ 2 & 3 & 1 \\ 5 & 1 & 1 \end{bmatrix}
=1(3-1)-0+2(2-15)=2-26=-24,\quad x_1=\frac{-24}{-12}=2.
\]
Replace column 2 by $\vec b$:
\[
\det\mat A_2=\det\begin{bmatrix} 1 & 1 & 2 \\ -1 & 2 & 1 \\ 2 & 5 & 1 \end{bmatrix}
=1(2-5)-1(-1-2)+2(-5-4)=-3+3-18=-18,\quad x_2=\frac{-18}{-12}=\tfrac32.
\]
Replace column 3 by $\vec b$:
\[
\det\mat A_3=\det\begin{bmatrix} 1 & 0 & 1 \\ -1 & 3 & 2 \\ 2 & 1 & 5 \end{bmatrix}
=1(15-2)-0+1(-1-6)=13-7=6,\quad x_3=\frac{6}{-12}=-\tfrac12.
\]
So $\vec x=\big(2,\tfrac32,-\tfrac12\big)$. Quick check on the first equation:
$1\cdot2+0\cdot\tfrac32+2\cdot(-\tfrac12)=2-1=1=b_1$. Good.

\begin{trapbox}
Cramer's rule needs $\det\mat A\neq0$. If $\det\mat A=0$ the formula divides by zero and is silent --
it cannot distinguish ``no solution'' from ``infinitely many.'' That is a real limitation: when a
system is singular, the determinant ratios collapse and you must fall back on elimination to learn
which case you are in. Cramer answers only the square, unique-solution question.
\end{trapbox}

\section{The Geometry of Cramer's Rule: Each Unknown Is a Ratio of Volumes}

Formulas are easier to trust when you can see them. Cramer's rule has a direct geometric reading,
and it is the reason this rule belongs in a geometry-first book at all.

Recall the picture of $\mat A\vec x=\vec b$ as a statement about columns. Writing $\mat
A=[\,\vec a_1\;\vec a_2\;\cdots\;\vec a_n\,]$, the equation says
\[
x_1\vec a_1+x_2\vec a_2+\cdots+x_n\vec a_n=\vec b.
\]
The unknowns $x_i$ are the \emph{coordinates of $\vec b$ in the basis of columns}. Cramer's rule
recovers each coordinate as a volume ratio. Let us see exactly why in $\R^2$, where ``volume'' means
area.

Take columns $\vec a_1,\vec a_2$ in the plane. The determinant $\det\mat A$ is the signed area of the
parallelogram they span. Now $\vec b=x_1\vec a_1+x_2\vec a_2$. Consider the parallelogram built from
$\vec b$ and $\vec a_2$ -- its area is $\det[\,\vec b\;\;\vec a_2\,]=\det\mat A_1(\vec b)$. Replace
$\vec b$ by its expansion:
\[
\det[\,\vec b\;\;\vec a_2\,]=\det[\,x_1\vec a_1+x_2\vec a_2\;\;\vec a_2\,]
= x_1\det[\,\vec a_1\;\;\vec a_2\,]+x_2\det[\,\vec a_2\;\;\vec a_2\,].
\]
The second determinant is zero (two equal columns). So $\det[\,\vec b\;\;\vec a_2\,]=x_1\det\mat A$,
giving $x_1=\det\mat A_1(\vec b)/\det\mat A$ -- Cramer's rule, derived from pure area.

\begin{intuitionbox}
Here is the one-sentence geometry. \textbf{$x_i$ is the factor by which you must stretch the $i$-th
column so that, together with the \emph{other} columns, it spans the same volume as $\vec b$ does with
those columns.} Replacing column $i$ by $\vec b$ and taking the area ratio peels off exactly that
stretch factor, because the other columns are shared and the determinant is linear in the column you
changed.
\end{intuitionbox}

\begin{figure}[H]
\centering
\begin{tikzpicture}[scale=1.05, >={Stealth[length=2.4mm]}]
  \def\ax{2.6}\def\ay{0.4}   % a1
  \def\bx{0.7}\def\by{2.1}   % a2
  % b = 0.8 a1 + 0.9 a2
  \pgfmathsetmacro{\rx}{0.8*\ax+0.9*\bx}
  \pgfmathsetmacro{\ry}{0.8*\ay+0.9*\by}
  % axes
  \draw[->, spacegray] (-0.3,0) -- (4.2,0) node[right, font=\footnotesize] {$x$};
  \draw[->, spacegray] (0,-0.3) -- (0,3.6) node[above, font=\footnotesize] {$y$};
  % full parallelogram a1,a2
  \fill[sourceblue!14] (0,0) -- (\ax,\ay) -- (\ax+\bx,\ay+\by) -- (\bx,\by) -- cycle;
  \draw[sourceblue!60, dashed] (0,0) -- (\ax,\ay) -- (\ax+\bx,\ay+\by) -- (\bx,\by) -- cycle;
  % parallelogram b, a2 (the A_1(b) area)
  \fill[vecred!16] (0,0) -- (\rx,\ry) -- (\rx+\bx,\ry+\by) -- (\bx,\by) -- cycle;
  \draw[vecred!70, thick] (0,0) -- (\rx,\ry) -- (\rx+\bx,\ry+\by) -- (\bx,\by) -- cycle;
  % decomposition line: sliding along a2 direction
  \draw[quotientorange, thick, dashed] (\rx,\ry) -- (0.8*\ax, 0.8*\ay)
    node[pos=0.45, right=2pt, font=\footnotesize, fill=white, inner sep=1pt] {shear $\parallel \vec a_2$};
  \fill[quotientorange] (0.8*\ax, 0.8*\ay) circle (1.5pt)
    node[above=2pt, font=\footnotesize, quotientorange!90!black] {$x_1\vec a_1$};
  % vectors
  \draw[->,vecblue,very thick] (0,0) -- (\ax,\ay) node[right=2pt]{$\vec a_1$};
  \draw[->,targetgreen,very thick] (0,0) -- (\bx,\by) node[left]{$\vec a_2$};
  \draw[->,vecred,very thick] (0,0) -- (\rx,\ry) node[above right=1pt]{$\vec b$};
  % Area labels with generous backgrounds
  \node[sourceblue!90!black, fill=white, inner sep=1.5pt, rounded corners=2pt, font=\footnotesize\bfseries]
    at (1.3,0.95) {$\text{Area} = \det\mat A$};
  \node[vecred!90!black, fill=white, inner sep=1.5pt, rounded corners=2pt, font=\footnotesize\bfseries]
    at (1.8,3.05) {$\text{Area} = \det\mat A_1(\vec b)$};
\end{tikzpicture}
\caption{Cramer's rule in the plane. The blue parallelogram on $\vec a_1,\vec a_2$ has signed area $\det\mat
A$. The red parallelogram on $\vec b,\vec a_2$ has area $\det\mat A_1(\vec b)$. Because $\vec b = x_1\vec a_1 + x_2\vec a_2$,
replacing $\vec a_1$ with $\vec b$ scales the base by $x_1$ and then shears parallel to $\vec a_2$ (orange dashed line).
By Cavalieri's principle, shearing parallel to $\vec a_2$ contributes zero extra area against $\vec a_2$, so the red
area is exactly $x_1$ times the blue area: $x_1 = \det\mat A_1(\vec b)/\det\mat A$.}
\label{ch22:fig-cramer}
\end{figure}

\begin{walkbox}
Consider why the component of $\vec b$ along $\vec a_2$ contributes nothing to the area
$\det[\,\vec b\;\vec a_2\,]$. A parallelogram whose edges are collinear collapses to zero
area. Consequently, the determinant $\det[\,\vec b\;\vec a_2\,]$ isolates the component of
$\vec b$ in the direction of $\vec a_1$, which is precisely $x_1\vec a_1$. This geometric
cancellation explains why Cramer's rule computes each coordinate individually.
\end{walkbox}

\begin{aestheticbox}
The fundamental reason Cramer's rule exists is that the determinant is \emph{multilinear} -- linear
in each column separately -- and \emph{alternating} -- zero when two columns coincide. Those two
properties alone force the volume ratio to extract a single coordinate. You will meet the same pair
of properties in exterior algebra and differential forms: an alternating multilinear form naturally
measures oriented volume.
\end{aestheticbox}

\section{When These Formulas Help, and When Elimination Wins}\label{ch22:sec-when}

Here is the verdict, laid out plainly.

\textbf{Elimination wins for numbers.} To solve $\mat A\vec x=\vec b$ for a specific numeric $\mat A$
and $\vec b$, Gaussian elimination costs about $n^3/3$ operations and solves all $n$ unknowns at once.
Cramer's rule asks you to compute $n+1$ separate determinants; even computing each determinant
\emph{by elimination} costs $\sim n^3/3$, so Cramer is roughly $n+1$ times more expensive -- and if
you compute those determinants by cofactor expansion the cost explodes to $\sim n\cdot n!$. For
anything beyond about $3\times3$, Cramer and cofactors are computationally hopeless. They are also
numerically fragile: dividing by a tiny determinant amplifies round-off error badly.

\textbf{Cofactors and Cramer win for structure, symbols, and small sizes.} Reach for them when:
\begin{itemize}
\item The matrix is \textbf{$2\times2$ or $3\times3$} and you want a quick closed form -- the $2\times
  2$ inverse recipe is the adjugate, and you should keep it at your fingertips.
\item You need a \textbf{single unknown}, not the whole solution. Cramer gives $x_3$ alone without
  ever finding $x_1,x_2$ -- occasionally handy in a derivation or an exam.
\item You are working \textbf{symbolically} -- the entries are letters or parameters -- and you want
  to see \emph{how} the solution depends on them. The explicit polynomial-over-determinant form is
  exactly right for sensitivity analysis and for proofs.
\item You are \textbf{proving a theorem}. The adjugate identity, the continuity of $\mat A\inv$ in its
  entries, the integrality results in the connection box -- all rest on these formulas, not on
  elimination.
\end{itemize}

\begin{focusbox}
Rule of thumb: \textbf{to get a number, eliminate; to get a formula or a proof, use cofactors.}
Cramer's rule and the adjugate are the determinant's \emph{theoretical} face. They tell you what the
answer \emph{is} as an expression, which is a different and sometimes more valuable thing than telling
you what it equals as a decimal.
\end{focusbox}

\begin{applicationbox}
\textbf{Where the adjugate quietly shows up in modern work.} In computer graphics and physics
engines, surface \emph{normals} must be transformed correctly when an object is deformed by a matrix
$\mat M$. A normal does not transform by $\mat M$ -- it transforms by the \emph{inverse transpose}
$(\mat M\inv)\T$, and for the common case where you only care about direction, this is proportional to
the \emph{transpose of the adjugate}, $\operatorname{adj}(\mat M)\T$. Because the adjugate avoids the
$1/\det$ division, graphics code often uses it directly: it stays finite and correct even when $\det
\mat M$ is small, and you only normalize the resulting direction at the end. Meanwhile in statistics,
Cramer's rule is the cleanest way to \emph{see} how a single regression coefficient depends on the
data: each coefficient in the normal equations is one determinant ratio, which makes the effect of
adding or perturbing a variable transparent -- even though the coefficients themselves are computed by
the far more stable $\mat Q\mat R$ or Cholesky methods, never by literally evaluating Cramer on a
large design matrix. The formula is the explanation; elimination is the calculator.
\end{applicationbox}

\begin{connectionbox}
\textbf{Cramer's rule and the implicit function theorem.} In multivariable calculus you ask: if
several quantities are tied together by equations, how does one of them change when you nudge an
input? Linearizing the equations turns that question into a linear system whose matrix is the
\emph{Jacobian} -- the matrix that collects all first-order partial derivatives, acting as the local best linear approximation. The implicit function theorem says you can solve for the derivative exactly when the
Jacobian determinant is nonzero -- and the derivative you get is, term by term, a Cramer's-rule ratio
of determinants. So the condition ``$\det\neq0$'' that lets Cramer run is the very same condition that
lets calculus invert a map locally. The two subjects are telling the same story: a nonzero determinant
means the linear approximation is invertible, so the answer is a clean ratio of volumes.
\end{connectionbox}

\begin{historybox}
In 1750, the Swiss mathematician Gabriel Cramer published an influential treatise on algebraic curves with a very concrete geometric puzzle in mind: finding the equation of a plane curve passing through a specified set of points. Fitting a conic section---an ellipse, parabola, or hyperbola---through five points required solving five simultaneous linear equations in five unknowns. For cubic curves, nine equations were needed. Cramer ran straight into a famous conundrum known as Cramer's paradox: two cubic curves can intersect in nine points, yet nine points were supposed to uniquely pin down a single cubic. To understand how systems of equations could become redundant or singular, Cramer sought a general algebraic formula for the unknowns.

At the time, mathematicians had no matrix notation, no vectors, no systematic elimination tables, and no subscript indices ($a_{ij}$). They wrote systems out with alphabetic letters ($a, b, c, d, \dots$) and untangled them through pages of exhausting substitution. In an appendix to his 1750 book, Cramer introduced an ingenious index notation that allowed him to state his celebrated rule: each unknown in an $n\times n$ system is the ratio of two determinants. 

Yet this history holds a famous priority controversy. Two years earlier, in 1748, the Scottish mathematician Colin Maclaurin published the exact same determinant formulas for $2\times2$ and $3\times3$ systems in his posthumous \emph{Treatise of Algebra} (having discovered them as early as 1729). Maclaurin published first, but Cramer's superior notation for arbitrary $n$ captured Europe's attention, and history permanently attached Cramer's name to the rule. Two decades later, Pierre-Simon Laplace (1772) showed how to peel any determinant apart by expanding it recursively along any row or column. In 1812, Augustin-Louis Cauchy synthesized the theory, established the word ``determinant'' in its modern sense, and proved the adjugate identity.

A common misconception is that Cramer's rule serves as a practical numerical method for solving systems on a computer. In practice, computing Cramer's rule via cofactor expansion requires an impractical number of operations. Solving an $n\times n$ system requires $n+1$ determinants. Peeling each into sub-determinants takes $\sim n!$ steps, demanding $(n+1)!$ arithmetic operations in total. For a modest $100\times 100$ system, $(101)! \approx 9.3 \times 10^{157}$ operations---vastly more than the number of seconds since the Big Bang ($\approx 4.35 \times 10^{17}$). Meanwhile, Gaussian elimination solves that same system in less than a millisecond.

If Cramer's rule is impractical for large-scale computation, what did it unlock? It gave mathematics an exact symbolic formula for matrix inverses and solutions, proving that solutions depend smoothly and differentiably on matrix entries---the cornerstone of the implicit function theorem in multivariable calculus and sensitivity analysis in economics.
\end{historybox}

\section{Looking Back}

We took the single number from the previous chapter -- the determinant as signed volume -- and learned
to take it apart and put it back together.

\begin{itemize}
\item \textbf{Cofactor expansion} (Theorem~\ref{ch22:thm-expansion}) computes a determinant along any
  row or column as a signed sum of smaller determinants. Expand where the zeros are. It is recursive,
  exact, and -- for large matrices -- factorially slow.
\item \textbf{Special matrices} fall out instantly: a triangular determinant is the product of its
  diagonal (Theorem~\ref{ch22:thm-triangular}), which is exactly why elimination computes determinants
  cheaply. And $\det(c\mat A)=c^n\det\mat A$, not $c\det\mat A$.
\item \textbf{The adjugate} (Theorem~\ref{ch22:thm-adj}) packages the cofactors into an explicit
  inverse $\mat A\inv=\frac{1}{\det\mat A}\operatorname{adj}\mat A$. Its off-diagonal entries vanish
  because pairing one row's entries with another row's cofactors evaluates a determinant with
  a repeated row.
\item \textbf{Cramer's rule} (Theorem~\ref{ch22:thm-cramer}) reads each unknown
  $x_i=\det\mat A_i(\vec b)/\det\mat A$ as a ratio of determinants, and geometrically as a ratio of
  volumes (Figure~\ref{ch22:fig-cramer}) -- the determinant's multilinear, alternating nature
  extracting one coordinate at a time.
\item \textbf{The verdict:} these are tools of \emph{insight}, not of \emph{computation}. To
  get a number, eliminate. To get a formula or a proof, reach for cofactors.
\end{itemize}

In the next chapters we apply the determinant to characteristic equations. The equation
$\det(\mat A-\lambda\mat I)=0$ defines the \emph{characteristic polynomial}, whose roots are
the eigenvalues of $\mat A$. Cofactor expansion is how you write that polynomial down by hand for
a $2\times2$ or $3\times3$ matrix, so the mechanics practiced here carry directly forward.

\section{Exercises}

\begin{exercise}[Mental check: determinant rules]\label{ch22:ex-mental-scaling}
A $4\times4$ matrix $\mat A$ has $\det\mat A = 3$. Without writing on paper, determine:
\begin{enumerate}[label=(\alph*)]
\item $\det(2\mat A)$.
\item $\det(-\mat A)$.
\item The determinant of the matrix obtained by swapping rows 1 and 3, and then multiplying row 2 by $-5$.
\item $\det(\mat A\inv)$.
\end{enumerate}
\end{exercise}

\begin{exercise}[Mental check: cofactors and duplicated rows]\label{ch22:ex-mental-repeated}
Suppose $\mat A$ is a $5\times5$ matrix whose second and fourth rows are identical.
\begin{enumerate}[label=(\alph*)]
\item What is $\det\mat A$?
\item Without doing any arithmetic, what is the matrix product $\mat A\,\operatorname{adj}(\mat A)$?
\item Look at entry $(2,4)$ of $\mat A\,\operatorname{adj}(\mat A)$. What combination of entries and cofactors forms this entry, and why does it vanish?
\end{enumerate}
\end{exercise}

\begin{exercise}\label{ch22:ex-warmup-expand}
Compute $\det\begin{bmatrix} 3 & 1 & 0 \\ 2 & -1 & 4 \\ 0 & 5 & 2\end{bmatrix}$ twice: once by
expanding along the first row, once by expanding along the first column. Confirm you get the same
number, and say which expansion required less arithmetic and why.
\end{exercise}

\begin{exercise}\label{ch22:ex-zeros}
Find a row or column with the most zeros and use it to evaluate
$\det\begin{bmatrix} 4 & 0 & 0 & 1 \\ 7 & 2 & 0 & 0 \\ 1 & 3 & -1 & 0 \\ 0 & 6 & 0 & 2\end{bmatrix}$.
(Hint: the third column has three zeros.)
\end{exercise}

\begin{exercise}\label{ch22:ex-triangular}
Without expanding, write down $\det\begin{bmatrix} 5 & 99 & -3 & 7 \\ 0 & 2 & 11 & 8 \\ 0 & 0 & -1 & 4
\\ 0 & 0 & 0 & 3\end{bmatrix}$. Then explain in one sentence why all the entries above the diagonal
were irrelevant.
\end{exercise}

\begin{exercise}\label{ch22:ex-scaling}
Let $\mat A$ be a $3\times3$ matrix with $\det\mat A=4$. Compute $\det(2\mat A)$, $\det(\mat A\T)$,
$\det(\mat A\inv)$, and $\det(\mat A^2)$. State the general rule you used in each case.
\end{exercise}

\begin{exercise}\label{ch22:ex-adj2x2}
For $\mat A=\begin{bmatrix} 4 & 3 \\ 6 & 5\end{bmatrix}$, write down the cofactor matrix, the adjugate,
and $\mat A\inv$. Verify $\mat A\,\operatorname{adj}\mat A=(\det\mat A)\mat I$ by direct multiplication.
\end{exercise}

\begin{exercise}\label{ch22:ex-cramer2}
Use Cramer's rule to solve $\begin{bmatrix} 3 & 2 \\ 1 & 4\end{bmatrix}\vec x=\begin{bmatrix} 7 \\
9\end{bmatrix}$. Then check your answer by substitution.
\end{exercise}

\begin{exercise}\label{ch22:ex-cramer-one}
For the system $\begin{bmatrix} 1 & 2 & -1 \\ 0 & 3 & 1 \\ 2 & 0 & 1\end{bmatrix}\vec
x=\begin{bmatrix} 3 \\ 5 \\ 4\end{bmatrix}$, use Cramer's rule to find \emph{only} $x_2$, without
computing $x_1$ or $x_3$. (This illustrates the practical advantage of Cramer's rule: isolating a single unknown without solving the entire system.)
\end{exercise}

\begin{exercise}\label{ch22:ex-adj3x3}
Compute $\mat A\inv$ for $\mat A=\begin{bmatrix} 2 & 0 & 1 \\ 1 & 1 & 0 \\ 0 & 1 & 3\end{bmatrix}$ by
the adjugate method: find $\det\mat A$, all nine cofactors, the adjugate, and divide. Then verify one
column of $\mat A\mat A\inv=\mat I$.
\end{exercise}

\begin{exercise}\label{ch22:ex-cost}
A computer does $10^9$ multiplications per second. Estimate, in seconds, how long a cofactor expansion
($\sim n!$ multiplications) would take for $n=15$ versus how long elimination ($\sim n^3/3$) would
take. Use your estimates to argue, in a sentence or two, why no serious software computes determinants
by cofactors.
\end{exercise}

\begin{exercise}\label{ch22:ex-geom-cramer}
In $\R^2$ let $\vec a_1=\colvec{3\\0}$, $\vec a_2=\colvec{1\\2}$, and $\vec b=\colvec{5\\4}$. Compute
$x_1,x_2$ by Cramer's rule, then verify geometrically: confirm that the area of the parallelogram on
$\vec b,\vec a_2$ equals $x_1$ times the area of the parallelogram on $\vec a_1,\vec a_2$. Sketch both
parallelograms.
\end{exercise}

\begin{exercise}\label{ch22:ex-proof-singular}
\textbf{(Proof.)} Use the adjugate identity $\mat A\,\operatorname{adj}\mat A=(\det\mat A)\mat I$ to
prove that if $\det\mat A=0$ then $\mat A\,\operatorname{adj}\mat A=\mat 0$. Then argue that if in
addition $\operatorname{adj}\mat A\neq\mat 0$, the matrix $\mat A$ must be singular with a nontrivial
null space, and that the nonzero columns of $\operatorname{adj}\mat A$ lie in that null space.
\end{exercise}

\begin{exercise}\label{ch22:ex-proof-detadj}
\textbf{(Challenge.)} For an invertible $n\times n$ matrix $\mat A$, take determinants of both sides
of $\mat A\,\operatorname{adj}\mat A=(\det\mat A)\mat I$ to show that
$\det\big(\operatorname{adj}\mat A\big)=(\det\mat A)^{\,n-1}$. Check your formula on a $2\times2$ and a
$3\times3$ example.
\end{exercise}
